

\documentclass[12pt]{amsart}

\usepackage[OT2, T1]{fontenc}

\usepackage{amscd}
\usepackage{amsmath}
\usepackage{amssymb}
\usepackage{mathrsfs} 			% for \mathscr (script letters)
\usepackage{units}
\usepackage[all]{xy}

\usepackage{algorithm}
\usepackage{algorithmic}

\renewcommand\algorithmicrequire{\textbf{Input:}}
\renewcommand\algorithmicensure{\textbf{Output:}}

%
%------    GENERAL MACROS    -----
%
% Standard rings and fields, affine and projective space
%
\def\NZQ{\Bbb}               % the font for N,Z,Q,R,C
\def\NN{{\NZQ N}}
\def\QQ{{\NZQ Q}}
\def\ZZ{{\NZQ Z}}
\def\RR{{\NZQ R}}
\def\CC{{\NZQ C}}
\def\AA{{\NZQ A}}
\def\PP{{\NZQ P}}
\def\FF{{\NZQ F}}
\def\GG{{\NZQ G}}
\def\HH{{\NZQ H}}
%
%------------------------------------------------
% Symbols in "Fraktur"

\def\frk{\frak}               % font for "Fraktur"
\def\aa{{\frk a}}
\def\pp{{\frk p}}
\def\Pp{{\frk P}}
\def\qq{{\frk q}}
\def\Qq{{\frk Q}}
\def\mm{{\frk m}}
\def\Mm{{\frk M}}
\def\Phi{{\frk n}}
\def\Phi{{\frk N}}
%
%------------------------------------------------
% Small letters in bold
%
\def\ab{{\bold a}}
\def\bb{{\bold b}}
\def\xb{{\bold x}}
\def\yb{{\bold y}}
\def\zb{{\bold z}}
%
\def\opn#1#2{\def#1{\operatorname{#2}}} % to make operators
%------------------------------------------------
% Numerical invariants of rings, ideals, and modules
%
\opn\chara{char} \opn\length{\ell} \opn\pd{pd} \opn\rk{rk}
\opn\projdim{proj\,dim} \opn\injdim{inj\,dim} \opn\rank{rank}
\opn\depth{depth} \opn\sdepth{sdepth} \opn\fdepth{fdepth}
\opn\grade{grade} \opn\height{height} \opn\embdim{emb\,dim}
\opn\codim{codim}  \opn\min{min} \opn\max{max}
\def\OO{{\mathcal O}}
\opn\Tr{Tr} \opn\bigrank{big\,rank}
\opn\superheight{superheight}\opn\lcm{lcm}
\opn\trdeg{tr\,deg}%\emph{
\opn\reg{reg} \opn\lreg{lreg} \opn\ini{in} \opn\lpd{lpd}
\opn\size{size}
%------------------------------------------------
% Divisors
%
\opn\div{div} \opn\Div{Div} \opn\cl{cl} \opn\Cl{Cl}
%
%------------------------------------------------
% Subsets of the spectrum of a ring
%
\opn\Spec{Spec} \opn\Supp{Supp} \opn\supp{supp} \opn\Sing{Sing}
\opn\Ass{Ass} \opn\Min{Min}
%
%------------------------------------------------
% Standard operations on ideals and modules
%
\opn\Ann{Ann} \opn\Rad{Rad} \opn\Soc{Soc}
%
%------------------------------------------------
% Linear algebra and homology, endo- and automorphisms
%
\opn\Im{Im} \opn\Ker{Ker} \opn\Coker{Coker} \opn\Am{Am}
\opn\Hom{Hom} \opn\Tor{Tor} \opn\Ext{Ext} \opn\End{End}
\opn\Aut{Aut} \opn\id{id}  \opn\deg{deg}
\def\Frob{{\mathcal F}}
\opn\nat{nat}
\opn\pff{pf}%   \pf exists already
\opn\Pf{Pf} \opn\GL{GL} \opn\SL{SL} \opn\mod{mod} \opn\ord{ord}
\opn\Gin{Gin} \opn\Hilb{Hilb}
%
%------------------------------------------------
% Convexity
%
\opn\aff{aff} \opn\con{conv} \opn\relint{relint} \opn\st{st}
\opn\lk{lk} \opn\cn{cn} \opn\core{core} \opn\vol{vol}
\opn\link{link} \opn\star{star}
%------------------------------------------------
% Graded rings and Rees algebras
\opn\gr{gr}
\def\Rees{{\mathcal R}}
%
%------------------------------------------------
% Polynomials and power series
%
\def\poly#1#2#3{#1[#2_1,\dots,#2_{#3}]}
\def\pot#1#2{#1[\kern-0.28ex[#2]\kern-0.28ex]}
\def\Pot#1#2#3{\pot{#1}{#2_1,\dots,#2_{#3}}}
\def\konv#1#2{#1\langle #2\rangle}
\def\Konv#1#2#3{\konv{#1}{#2_1,\dots,#2_{#3}}}
%
%------------------------------------------------
% Direct and inverse limits
%
\opn\dirlim{\underrightarrow{\lim}}
\opn\inivlim{\underleftarrow{\lim}}
%
%
% Names with a meaning
%
\let\union=\cup
\let\sect=\cap
\let\dirsum=\oplus
\let\tensor=\otimes
\let\iso=\cong
\let\Union=\bigcup
\let\Sect=\bigcap
\let\Dirsum=\bigoplus
\let\Tensor=\bigotimes
%
%------------------------------------------------
%
\let\to=\rightarrow
\let\To=\longrightarrow
\def\Implies{\ifmmode\Longrightarrow \else
        \unskip${}\Longrightarrow{}$\ignorespaces\fi}
\def\implies{\ifmmode\Rightarrow \else
        \unskip${}\Rightarrow{}$\ignorespaces\fi}
\def\iff{\ifmmode\Longleftrightarrow \else
        \unskip${}\Longleftrightarrow{}$\ignorespaces\fi}
\let\gets=\leftarrow
\let\Gets=\longleftarrow
\let\followsfrom=\Leftarrow
\let\Followsfrom=\Longleftarrow
\let\:=\colon
%
%
%
\newtheorem{Theorem}{Theorem}[]
\newtheorem{Lemma}[Theorem]{Lemma}
\newtheorem{Corollary}[Theorem]{Corollary}
\newtheorem{Proposition}[Theorem]{Proposition}
\newtheorem{Examples}[Theorem]{Examples}
\newtheorem{Problem}[Theorem]{}
\newtheorem{Conjecture}[Theorem]{Conjecture}
\newtheorem{Question}[Theorem]{Question}
\newtheorem{variant}[Theorem]{Variant}

\theoremstyle{definition}
\newtheorem{Example}[Theorem]{Example}
\newtheorem{Definition}[Theorem]{Definition}
\newtheorem{Remark}[Theorem]{Remark}
\newtheorem{Remarks}[Theorem]{Remarks}


%------------------------The bpp stuff
\newtheoremstyle{subsection-tweak}
   {11pt}
   {3pt}%
   {}
   {}%
   {\bfseries}
   {}%
   {.5em}
   {\thmnumber{\@{#1}{}\@{#2}.}%
    \thmnote{~{\bfseries#3.}}}    

\newcounter{numberingbase}
\def\thenumberingbase{A.\arabic{numberingbase}}    

\theoremstyle{subsection-tweak}
\newtheorem{bpp}[Theorem]{}
\newtheorem{bppt}[numberingbase]{}
\newcommand{\bbpp}{\begin{bpp}}
\newcommand{\eepp}{\end{bpp}}
\newcommand{\bbppt}{\begin{bppt}}
\newcommand{\eeppt}{\end{bppt}}


\theoremstyle{theorem}
\newtheorem{TheoremT}[numberingbase]{Theorem}
\newtheorem{LemmaT}[numberingbase]{Lemma}
\newtheorem{CorollaryT}[numberingbase]{Corollary}
\newtheorem{PropositionT}[numberingbase]{Proposition}
\newtheorem{ExamplesT}[numberingbase]{Examples}
\newtheorem{ProblemT}[numberingbase]{}
\newtheorem{ConjectureT}[numberingbase]{Conjecture}
\newtheorem{QuestionT}[numberingbase]{Question}
\newtheorem{variantT}[numberingbase]{Variant}

\theoremstyle{definition}
\newtheorem{ExampleT}[numberingbase]{Example}
\newtheorem{DefinitionT}[numberingbase]{Definition}
\newtheorem{RemarkT}[numberingbase]{Remark}
\newtheorem{RemarksT}[numberingbase]{Remarks}




%%%%%%%%%%%%%%%%%%%%%%%%%

\newcommand{\val}{\mathrm{val}}		% Valuation

\newcommand{\x}{\text}
\newcommand{\rv}{\revise}
\newcommand{\rd}{\ready}
\newcommand{\q}{\quad}
%\newcommand{\qq}{\quad\quad}
\newcommand{\qqq}{\quad\quad\quad}
\newcommand{\qqqq}{\quad\quad\quad\quad}
\providecommand{\qxq}[1]{\quad\text{#1}\quad}
\providecommand{\qx}[1]{\quad\text{#1}}
\providecommand{\xq}[1]{\text{#1}\quad}
\newcommand{\tst}{\textstyle}
\newcommand*{\QED}{\hfill\ensuremath{\qed}}
\newcommand{\sU}{\mathscr{U}}
\newcommand{\msk}{\medskip}
\newcommand{\bsk}{\bigskip}
\DeclareMathOperator{\Frac}{Frac}		% Field of fractions
\newcommand{\wt}{\widetilde}
\newcommand{\ssk}{\smallskip}
\providecommand{\up}[1]{{\upshape(}#1{\upshape)}}
\providecommand{\ucolon}{{\upshape:} }
\newcommand{\surjects}{\twoheadrightarrow}
\newcommand{\hra}{\hookrightarrow}



% We like the var forms of some greek letters (as taught in German schools)
%
\let\epsilon\varepsilon
\let\phi=\varphi
%\let\kappa=\varkappa
%
%           We print on A4 paper
%
\textwidth=15cm \textheight=22cm \topmargin=0.5cm
\oddsidemargin=0.5cm \evensidemargin=0.5cm \pagestyle{plain}
%
%           The pf environment of AMSART needs a little help
%
\def\qed{\ifhmode\textqed\fi
      \ifmmode\ifinner\quad\qedsymbol\else\dispqed\fi\fi}
\def\textqed{\unskip\nobreak\penalty50
       \hskip2em\hbox{}\nobreak\hfil\qedsymbol
       \parfillskip=0pt \finalhyphendemerits=0}
\def\dispqed{\rlap{\qquad\qedsymbol}}
\def\noqed{\def\qed{\relax}}
%
% ------    END OF GENERAL MACROS    -------
\opn\dis{dis}
\def\pnt{{\raise0.5mm\hbox{\large\bf.}}}
\def\lpnt{{\hbox{\large\bf.}}}
\opn\Lex{Lex}

%-- macro for local cohomology-----------------------------
\def\Coh#1#2{H_{\mm}^{#1}(#2)}
%-- macro for a complicated condition for the extended
%-- Hochster's formula
\def\hchst#1{for all $u\in G(I)$ there  exists $i\notin {#1}$
such that $\nu_i(u) > a_i \geq 0$}

\begin{document}


\title{A special form of Zariski's Uniformization Theorem in positive characteristic.}

\author{ Dorin Popescu}

\address{Simion Stoilow Institute of Mathematics of the Romanian Academy,
Research unit 5, P.O. Box 1-764, Bucharest 014700, Romania,}

\address{University of Bucharest, Faculty of Mathematics and Computer Science
Str. Academiei 14, Bucharest 1, RO-010014, Romania,}

\address{ Email: {\sf dorin.m.popescu@gmail.com}}

\dedicatory{To the $80^{th}$ birthday   of Vasile Br\^ inz\u anescu.}


\begin{abstract} This is mainly a small exposition on extensions of valuation rings
as a filtered union of smooth algebras. \\
\vskip 0.3 cm
{\it Key words }: immediate extensions, smooth morphisms, Henselian rings, complete intersection algebras, pseudo convergent sequences, pseudo limits.   \\
 {\it 2020 Mathematics Subject Classification: Primary 13F30, Secondary 13A18,13F20,13B40.}
\end{abstract}

\maketitle
\section{Introduction}

A local ring   $(A,\mm)$ is  Noetherian if its  ideals are finitely generated. The maximum of the length of the chains
$p_0
\subsetneq p_1\subsetneq\ldots \subsetneq p_d$
 of its prime ideals is the Krull {\em dimension} of $A$. The Krull dimension of $A$ corresponds to dimension of algebraic varieties in algebraic geometry.
Another invariant is the  {\em embedding dimension} of $A$ which is 
 $\dim_{A/\mm} \mm/\mm^2$  as a linear space.

Let $C$ be the curve in the real plane
  ${\bf R}^2$ given by the equation   $X^2-Y^3=0$. The origin of the plane  $(0,0)$ is a point of  $C$, which is not smooth  in the usual terminology.  The dimension of the curve $C$ is $1$ in geometry  and its associated local ring  $B=({\bf R}[X,Y]/(Y^3-X^2))_{(X,Y)}$ in its point $(0,0)$ has a unique  big chain of prime ideals   $(0)\subset (X,Y)B$, that is $\dim B=1$. Note that 
  $\dim_{\bf R} (X,Y)B/(X,Y)^2B=2$.
on the other hand, the local ring   
   $B'=({\bf R}[X,Y]/(Y^3-X^2))_{(X-1,Y-1)}$ associated to  $C$ in its point $(1,1)$ has also a unique  big chain $(0)\subset (X-1,Y-1)B'$, that is the Krull dimension of   $B'$ is still $1$. Applying  Taylor's formula  for  $f=X^2-Y^3$ we get 
$$f=2(X-1)+(X-1)^2-3(Y-1) -3(Y-1)^2-(Y-1)^3$$ and so in the linear space 
  $( X-1,Y-1)B'/( X-1,Y-1)^2B'$ over  $\bf R$ there exists a linear dependence 
 between the elements induced by  
    $X-1,Y-1$ and so  the embedding dimension of  $B'$ is $1$.
 In general we have 
 $\dim A\leq \dim_{A/\mm} \mm/\mm^2 $ for a Noetherian local ring $(A,\mm)$. 
Note that this inequality is strict for $B$ but equality for $B'$. 


 A Noetherian local ring $(A,\mm)$ is {\em regular} if the above inequality is an equality.
In geometry the local regular rings correspond to the smooth points of a variety. This is a reason that  $(1,1)$ is a smooth point of $C$ but $(0,0)$ not.
  
  
Consider the transformation 
   ${\bf R}\to C$ given by  $u\to (u^3,u^2)$ and we see that so we obtain a desingularization $\bf R$ of $C$, the real line being smooth.
  Algebraically this means to embed the domain  
    $A={\bf R}[X,Y]/(X^2-Y^3)$ in the subring $B$ generated by   $X/Y$ in the fraction field    $K$ of  $A$, which is finite over $A$ because     $(X/Y)^2-Y$ is $0$ in $K$. Note that $A$ and $B$ have the same fraction field. Actually, $B$ is the integral closure of $A$ in $K$ and we may use the integral closure to define the desingularization of curves even when we replace $\bf R$  
 by a field of positive characteristic.
 
    The desingularization of an algebraic variety of higher dimension was a hard problem of the last century solved by H.Hironaka in 1964 in the case of characteristic zero. His proof followed the 
Zariski Uniformization theorem \cite{Z}, which says that 
     a valuation ring V containing a field $K$ of characteristic zero is a filtered union of its regular  $K$-subalgebras. In positive characteristic such theorem is missing.
     
     An  $A$-algebra $D$ is  {\em
smooth} if $D$ is a localization of an $A$-algebra of type  $(A[Y]/(f))_M$, where  $f=(f_1,\ldots ,f_r) $  is a system of polynomials
in  $Y=(Y_1,\ldots,Y_N)$, $N\geq r$ over  $A$ and  $M$ is an $r\times r$-minor of $(\partial f/\partial Y)$. If $A$ is a regular local ring then $D$ is too. If $K$ is a field of characteristic   zero and $D$ is a $K$-algebra, regular local ring, essentially of finite type over $K$ then $D$ is smooth over $K$. 

The Zariski's Uniformization Theorem could have the following form
     
  
\begin{Theorem} \label{T1}  (Zariski \cite{Z}) A valuation ring containing a field $K$ of characteristic zero is a    a filtered union of its smooth $K$-subalgebras.
\end{Theorem}  
   
   
   Let $k$ be a field of characteristic $2$ and $K=k(X)$ the field of fractions of the polynomial ring $k[X]$. Then $E=K[Y]/(Y^2-X)$ is a field, so a regular local ring,
   but $E/K$ is not separable, so $ E $ is not a smooth $K$-algebra. Thus in positive characteristic a similar statement of the above theorem needs to say more then the fact that a valuation ring containing a field $K$ is a filtered union of its regular local $K$-subrings.
   
 However, B. Antieau, R. Datta \cite[Theorem 4.1.1]{AD} showed the following theorem. 
\begin{Theorem} (\cite{AD})\label{T2}
  Every perfect valuation ring of characteristic $p>0$ is a filtered union of its smooth ${\bf F}_p$-subalgebras. 
 \end{Theorem}
 
 Recently, we showed the following theorem.
 
   \begin{Theorem}(\cite[Corollary 5]{P3})\label{T3} Let $V$ be a  valuation ring containing its residue field $k$   with a $\bf Z$-free value group $\Gamma$ (for example when it is finitely generated) and $K$ its fraction field. Assume that  $K=k(x,y)$ for some  algebraically independent elements $x=(x_i)_{i\in I},y=(y_j)_{j\in J}$ over $k$ such that $\val(y)$ is a basis in $\Gamma$. 
 Then $V$ is a filtered  union of its smooth $k$-subalgebras. In particular, $V$ is a filtered  union of its  $k$-subalgebras, which are regular local rings, essentially of finite type over $k$.
\end{Theorem}

It is the purpose of this paper to explain this result together with some of its applications.

\section{Immediate extension of valuation rings}


An {\it{immediate extension}} of valuation rings is an extension inducing trivial extensions on residue fields and group value extensions.

Using \cite[Theorem 1,VI,(10.3)]{Bou} we got in \cite[Lemma 26 (1)]{P} the following result.
   \begin{Lemma}(\cite{P})\label{l} Let $V$ be a  valuation ring containing its residue field $k$    with a $\bf Z$-free value group $\Gamma$ (for example when it is finitely generated), $\val$ its valuation and $K$ its fraction field. Assume that  $K=k(x)$ for some   elements $x=(x_i)_{i\in I}$ over $k$ such that $\val(x)$ is a basis in $\Gamma$. 
 Then $V$ is a filtered  union of its  $k$-subalgebras, which are localizations of some polynomial rings over $k$ in some $x$, so some smooth $k$-algebras.
\end{Lemma}

Let $V$ be a valuation ring  with a $\bf Z$-free  value group  $\Gamma$,  containing its residue field $k$  and $x$ some elements of $V$ such that $\val(x)$ is a basis of $\Gamma$.
Our idea was to split in two parts  a proof of a possible extension of Zariski's Uniformization Theorem in positive characteristic. Firstly we noticed that $W=V\cap k(x)$ is a  a filtered  union of its smooth $k$-subalgebras by Lemma \ref{l}. Then we try to find cases when the immediate extension $V/W$ is  a  a filtered  union of its smooth $W$-subalgebras. Unfortunately, in general  this is not the case (see \cite[Example 3.1.3]{Po1},\cite{O},  \cite[Example 3.6]{P4}) and for some positive results we have to state different other necessary results.

The above idea was used already in \cite{P}, where we proved the following result.

\begin{Theorem} (\cite[Theorem 2]{P})
 If $V\subset V'$ is an immediate extension of valuation rings   containing $\bf Q$
then $V'$ is a filtered direct colimit of smooth $V$-algebras. 
\end{Theorem}


Let $V$ be a valuation ring, $\lambda$ be a fixed limit ordinal  and $v=\{v_i \}_{i < \lambda}$ a sequence of elements in $V$ indexed by the ordinals $i$ less than  $\lambda$. Then $v$ is called \emph{pseudo convergent} if 

$\val(v_{i} - v_{i''} ) < \val(v_{i'} - v_{i''} )     \ \ \mbox{for} \ \ i < i' < i'' < \lambda$
(see \cite{Kap}, \cite{Sch}).
A  \emph{pseudo limit} of $v$  is an element $w \in V$ with 

$ \val(w - v_{i}) < \val(w - v_{i'}) \ \ \mbox{(that is,} \ \ \val(w -  v_{i}) = \val(v_{i} - v_{i'}) \ \ \mbox{for} \ \ i < i' < \lambda$.

k
\begin{Lemma}(\cite[Proposition 16]{P3}) \label{l1}
Let $V \subset V'$ be an  immediate extension of valuation rings,  $K, K'$ the fraction fields of $V, V'$, $\mm$ the maximal ideal of $V$ and  $y\in K'$ an element  which is not in  $K$. Assume that  $y$ is  a pseudo limit of a pseudo convergent sequence $v=(v_j)_{1\leq j<\lambda}$ over $V$, which has no pseudo limit in $K$.
 Set $y_j=(y-v_j)/(v_{j+1}-v_j)$.
Then for every nonzero polynomial  $g =\sum_{e=0}^m a_e Y^e\in V[Y]$ such that  $a_e=0$ for any $e>o$ multiple of $p$ if char\ $K=p>0$,  and every ordinal $1\leq \nu<\lambda$,  there exist some  $\nu<t<\lambda$ and a polynomial $g_1\in V[Y_t]$ such that 
$g(y)=g_1(y_t)$
 and    $g_1=g_1(0)+c  Y_t+g_2,$
for some $c\in V\setminus \{0\}$ and  $g_2\in  c\mm Y_t^2 V[Y_t]$. 
\end{Lemma}

For the proof very important is that changing $Y\to (Y-v_j)/(v_{j+1}-v_j)$ in $g$  we may assume that the nonzero coefficients of $g-g(0)$ have different values for $j$ large enough. It is worth to note that $\val(g(y)) $ could be different from $\val(g(y_j))$ for all $j$ as shows \cite[Example 16]{P4}. 

\begin{Remark} \label{r}
If the characteristic of $V$ is $p>0$ then taking above $g\in V[Y^p]$ we see that for any $t$ the polynomial $g_1=g(Y_t)\in V[Y_t^p]$ and so the coefficient of $Y_t$ in $g_1$ is zero.
\end{Remark}

When $p>0$ and there exist  in $g$ above some nonzero coefficients of some $Y^n$ with $n>0$ multiple of $p$ then we get the following lemma.

\begin{Lemma}(\cite[Corollary 17]{P3})\label{l2}
Let $V \subset V'$ be an  immediate extension of valuation rings,  $K, K'$ the fraction fields of $V, V'$, $\mm$ the maximal ideal of $V$ and  $y\in K'$ an  element  which is not in  $K$. 
Assume that char\ $K=p>0$ and $y$ is  a pseudo limit of a pseudo convergent sequence $v=(v_j)_{1\leq j<\lambda}$ over $V$, which  has no pseudo limit in $K$.
 Set $y_j=(y-v_j)/(v_{j+1}-v_j)$.
Then for every nonzero  polynomial  $g\in V[Y]$  and every ordinal $1\leq \nu<\lambda$ one of the following statements holds.

\begin{enumerate}
\item  There exist some  $\nu<t<\lambda$ and a polynomial $g_1\in V[Y_t]$ such that 
$g(y)=g_1(y_t)$
 and    $g_1=g_1(0)+c  Y_t+g_2,$ 
for some $c\in V\setminus \{0\}$ and  $g_2\in  c\mm Y_t^2 V[Y_t]$.
\item 
There exist some  $\nu<t<\lambda$ and a polynomial $g_1\in V[Y_t]$ such that 
$y g(y)=g_1(y_t)$
 and    $g_1=g_1(0)+c  Y_t+g_2,$ 
for some $c\in V\setminus \{0\}$ and  $g_2\in  c\mm Y_t^2 V[Y_t]$.
\end{enumerate}
\end{Lemma}

These lemmas are applied  to Lemma \ref{l3} given below together with the next lemma which is in fact our key lemma.

\begin{Lemma}(\cite[Lemma 21]{P3}) \label{k}
Let $V \subset V'$ be an  immediate extension of valuation rings,  $K, K'$ the fraction fields of $V, V'$, $\mm, \mm'$ the maximal ideals of $V,V'$, 
$(y_e)_{1\leq e\leq  m}$ some elements of $V'$ and $g_e\in V[Y_1,\ldots,Y_m]$, $1\leq e<m$ some polynomials such that the determinant of $((\partial g_e/\partial Y_i)((y_{e'}))_{1\leq e< m, 1\leq i<m}$ is not in $\mm'$.
 Assume that $y_m$ is transcendental over $K$ and $g_e((y_{e'})_{e'})=0$ for all $1\leq e<m$.
 Then $V[y_1,\ldots,y_m]_{\mm'\cap V[y_1,\ldots,y_m]}$ is  a smooth $V$-subalgebra  of $V'$.
\end{Lemma}

\begin{Lemma}(\cite[Lemma 22]{P3}) \label{l3}
 Let  $ V'$ be an  immediate extension  of a valuation ring $V$ containing a field, $\mm, \mm'$ the maximal ideals of $V,V'$, $K\subset K'$ their fraction field extension and $y\in V'$ an unit, transcendental element over $K$. 
Let  $f\in V[Y]$ be a nonzero polynomial and some $d\in V\setminus \{0\}$ and an unit $z\in V'$ such that $f(y)=d z$. Then there exists 
 a smooth $V$-subalgebra  of $V'$ containing $y,z$.  
\end{Lemma}

The above lemma gives an  idea of the proof of our  desingularization (in fact smoothification).
By \cite[Theorem 1]{Kap}  $y$ is a pseudo limit of a pseudo convergent sequence $v=(v_j)_{j<\lambda}$, which 
 has no pseudo limits in $K$.  Set $y_j=(y-v_j)/(v_{j+1}-v_j)$.


 By Lemmas \ref{l1}, \ref{l2} there exist some  $j<\lambda$ and a polynomial $h\in V[Y_j]$ such that either 
$f(y)=h(y_j)$, or $y f(y)=h(y_j)$
 and  $h=h(0)+c  Y_j+h',$
for some   $c\in V\setminus \{0\}$ and $h'\in  c\mm Y_j^2 V[Y_j]$.   If $c|d$ then $c|h(0)$ and $(h-dZ)/c$, or $(h-dYZ)/c$ is in the kernel of the map $V[Y_j,Z]\to V'$,
$(Y_j,Z)\to (y_j,z)$ and  $V[y_j,z]_{\mm'\cap V[y_j,z]}$ is smooth over $V$ by Lemma \ref{k} containing  $y,z$. 




\section{A smoothification theorem}

Mainly we try to explain here the following theorem

\begin{Theorem}(\cite[Theorem 23]{P3}) \label{T4}
 Let  $ V'$ be an  immediate extension  of a valuation ring $V$ containing a field, $\mm, \mm'$ the maximal ideals of $V,V'$, $K\subset K'$ their fraction field extension and $y_0\in V'$ an unit, transcendental element over $K$ with $K'=K(y_0)$. 
  Then $V'$ is  a filtered  union of its  smooth  $V$-subalgebras.
\end{Theorem}

For the proof we need also some extensions of Lemmas \ref{l1}, \ref{l2} given for two variables. Their proofs are similar but  more complicated.

\begin{Lemma}(\cite[Proposition 18]{P3}) \label{l4}
Let $V \subset V'$ be an  immediate extension of valuation rings,  $K, K'$ the fraction fields of $V, V'$, $\mm$ the maximal ideal of $V$ and  $y_1,y_2\in K'$ two elements,  which are not in $K$. Assume that $y_i$, $i=1,2$ is  a pseudo limit of a pseudo convergent sequence $v_i=(v_{i,j})_{1\leq j<\lambda_i}$, $i=1,2$ over $V$, which   have no pseudo limit in $K$.
 Set $y_{ij}=(y_i-v_{i,j})/(v_{i,j+1}-v_{i,j})$.
Then for every nonzero polynomial   $g=\sum_{e_1,e_2=0}^m a_{e_1,e_2} Y_1^{e_1} Y_2^{e_2} \in V[Y_1,Y_2]$, such that $a_{e_1,e_2}=0$ if char\ $K=p>0$ and at least one of $e_1,e_2$ is a nonzero multiple of  $p$,  and every ordinals $1\leq \nu_i<\lambda_i$,  $i=1,2$  there exist some  $\nu_i<t_i<\lambda_i$, $i=1,2$ and a polynomial $g_1\in V[Y_{1,t_1},Y_{2,t_2}]$ such that 
$g(y_1,y_2)=g_1(y_{1,t_1},y_{2,t_2})$
 and  $g_1=g_1(0)+c_1  Y_{1,t_1}+c_2 Y_{2,t_2}+g_2,$
for some   $c_1,c_2\in V, $ at least one of them nonzero, and $g_2\in  (c_1,c_2)\mm (Y_{1,t_1},Y_{2,t_2})^2 V[Y_{1,t_1},Y_{2,t_2}]$.   
\end{Lemma}



\begin{Lemma}(\cite[Corollary 20]{P3}) \label{l5}
Let $V \subset V'$ be an  immediate extension of valuation rings,  $K, K'$ the fraction fields of $V, V'$, $\mm$ the maximal ideal of $V$ and  $y_1,y_2\in K'$ two  elements  which are not in  $K$. 
Assume that char\ $K=p>0$ and $y_i$, $i=1,2$ are   pseudo limits of two pseudo convergent sequences $v_i=(v_{i,j})_{1\leq j<\lambda_i}$, $i=1,2$ over $V$, which have no pseudo limits in $K$.
 Set $y_{i,j}=(y_i-v_{i,j})/(v_{i,j+1}-v_{i,j})$.
Then for every nonzero polynomial  $f\in V[Y_1,Y_2]$  and every two ordinals $\nu_i<\lambda_i$ 
there exist some  $1\leq \nu_i<t_i<\lambda_i$, $i=1,2$ and a polynomial $g\in V[Y_{1,t_1}, Y_{2,t_2}]$
 such that $g=g(0)+c_1  Y_{1,t_1}+c_2  Y_{2,t_2}+g',$ 
for some $c_1,c_2\in V$, at least one of them nonzero,   $g'\in  (c_1,c_2)\mm (Y_{1,t_1},Y_{2,t_2})^2   V[Y_{1,t_1}, Y_{2,t_2}]$ 
and $g(y_{1,t_1},y_{2,t_2})$ is one of the following elements
$f(y_1,y_2)$, \ $y_1f(y_1,y_2)$, \ $y_2f(y_1,y_2)$, \ $y_1y_2f(y_1,y_2)$.
\end{Lemma}
The proof of Theorem \ref{T4} goes by induction  using Lemmas \ref{l1}, \ref{l2},   \ref{k}, \ref{l4}, \ref{l5}
as  in the proof of Lemma \ref{l3}. If
 $f_i(y)=d_iz_i$, $i=1,\ldots, s$
for some $f_i\in V[Y]$, $d_i\in V\setminus \{0\}$, $z_i\in V'$ then there exists a smooth $V$-subalgebra of $V'$  containing $y,(z_i)$.



Using by induction Theorem \ref{T4} we can extend it  for the case when $V\subset V'$ is a pure transcendental extension.
\begin{Theorem}(\cite[Theorem 4]{P3})\label{T5} Let  $V\subset V'$ be an immediate  extension of valuation rings   containing  a field and $K\subset K'$ its fraction field extension. If $K'=K(x)$ for some algebraically independent system of elements $x$ over $K$ 
then $V'$ is a filtered union of smooth $V$-subalgebras of $V'$. 
\end{Theorem}

{\bf Proof of Theorem \ref{T3}} 

Use  Theorem \ref{T7} and Lemma \ref{l}.

\section{Some applications}

A {\em Henselian} local ring is a local ring 
 $(A,\mm)$ for which the Implicit Function Theorem hods, that is 
 for every system of polynomials  $f=(f_1,\ldots ,f_r) $ of $A[Y]$,
$Y=(Y_1,\ldots,Y_N)$, $N\geq r$ over  $A$ and every solution  $y=(y_1,\ldots,y_N)\in A^N$     modulo $\mm$ of $f$ in   $A$ such that an $r\times r$-minor $M$ of Jacobian matrix $(\partial f_i/\partial Y_j)_{1\leq i\leq r,1\leq j\leq N}$  satisfies $M(y)\not \in \mm$ there exists a solution $\tilde y$ of $f=0$ in  $A$ with $y\equiv {\tilde y}$ modulo $\mm$.
 Thus  $(A,\mm)$ is Henselian if for any smooth $A$-algebra $B$  any $A$-morphism $B\to A/\mm$ can be lifted to an $A$-morphism $w:B\to A$. Thus  $w$ is a retraction of $A\to B$.

   

As in \cite[Proposition 18]{P1}(see also \cite[Corollary 24]{P3}, \cite[Proposition 1]{P4}) we have the following corollary.

\begin{Corollary}(\cite{P4})\label{C} Let $V\subset V'$ be an immediate extension of valuation rings containing a field and $K\subset K'$ its fraction field extension. If $K'=K(x)$ for some algebraically independent system of elements $x$ over $K$ 
and $V$ is Henselian then every system of polynomials over $V$, which has a solution in $V' $ has also one in $V$.
\end{Corollary}

The ideas of this corollary were used to show some conjectures of M. Artin from \cite{Ar} in 
\cite[Theorems 1.3, 1.4]{Po0} (see also \cite[Theorem 5.3.1]{I}) using the so called the General N\'eron desingularization (see  \cite[Theorem 2.5]{Po0}, \cite[Theorem 1.1]{S} and \cite[Theorem 5.2.56]{I}).


 


 





For a general immediate extension of valuation rings $V\subset V'$ we should expect that $V'$ is a  filtered  union of its complete intersection $V$-subalgebras. 
 
A {\em complete intersection} $V$-algebra {\it{essentially of finite type}} is a local $V$-algebra of type $C/(P)$, where $C$ is a localization of a polynomial $V$-algebra of finite type and $P$ is a regular sequence of elements of $C$. 

\begin{Theorem}  (\cite[Theorem 1]{P2})\label{T6} Let  $ V'$ be an  immediate extension  of a valuation ring $V$  and  $K\subset K'$ the fraction field extension. If $K'/K$ is algebraic  then $V'$ is a filtered
 union of its complete intersection $V$-subalgebras of finite type.
\end{Theorem}

As a consequence of Theorems \ref{T5}, \ref{T6} the next theorem follows.


\begin{Theorem} (\cite[Theorem 6]{P3}) \label{T7} Let  $ V'$ be an  immediate extension  of a valuation ring $V$  and  $K\subset K'$ the fraction field extension. Then $V'$ is a filtered
 union of its complete intersection $V$-subalgebras of finite type.
\end{Theorem}

\begin{Corollary} (\cite[Corollary 7] {P3})Let $V$ be a  valuation ring containing its residue field $k$   with a $\bf Z$-free value group $\Gamma$. 
 Then $V$ is a filtered  union of its complete intersection $k$-subalgebras. 
\end{Corollary}
For the proof choose  some elements $y$ in $V$ such that $\val(y)$ is a basis in $\Gamma$  and apply Lemma \ref{l}   for $V\cap k(y) $ and the above theorem for the immediate extension $V\cap k(y)\subset V$.

 

\begin{thebibliography}{99}
\normalsize{



\bibitem{AD} B.\  Antieau, R.\ Datta, {\em Valuation rings are derived splinters}, Math. Zeitschrift, {\bf 299}, (2021, 827-851,  arxiv/AG:2002.010627v1.

\bibitem{Ar} M.\ Artin,  {\em Construction techniques for algebraic spaces}, Actes Congres Intern. Math. {\bf 1}, (1970), 419-423.




\bibitem{Bou}   N.\ Bourbaki,  {\em \'El\'ements de math\'ematique. Alg\`ebre commutative}, chap. I-VII, Hermann (1961, 1964,
1965); chap. VIII-X, Springer (2006, 2007) (French).


\bibitem{I} C. Ionescu,{\em Classes of Good Noetherian Rings}, Birkh\"auser, (2023).




\bibitem{Kap} I.\ Kaplansky, {\em Maximal fields with valuations}, Duke Math. J. {\bf 9} (1942), 303-321.







\bibitem{O} A.\  Ostrowski, {\em Untersuchungen zur arithmetischen Theorie der K\"orper}, Math. Z. {\bf39} (1935), no. 1,
321-404.

 \bibitem{Po0} D.\ Popescu, {\em General Neron Desingularization and approximation}, Nagoya Math. J., {\bf 104}, (1986), 85-115.
 
\bibitem{Po0'} D.\ Popescu, {\em Letter to the Editor, General Neron Desingularization and approximation}, Nagoya Math. J., {\bf 118}, (1990), 45-53.
 



 \bibitem{Po1} D.\ Popescu, {\em Algebraic extensions of valued fields}, J. Algebra {\bf 108}, (1987), no. 2,
513-533.

 \bibitem{P} D.\ Popescu, {\em N\'eron desingularization of extensions of valuation rings with an Appendix by K\k{e}stutis \v{C}esnavi\v{c}ius}, in   "Transcendence in Algebra, Combinatorics, Geometry and Number Theory",  Eds. Alin Bostan, Kilian
Raschel, Springer Proceedings in Mathematics and Statistics 373,(2021), 275-307, 
   arxiv/AC:1910.09123v4.
   
\bibitem{P1}  D.\ Popescu, {\em Extensions of valuation rings containing $\bf Q$ as limits of smooth algebras}, Bull. Math.  Soc. Sci. Math. Roumanie, {\bf 65(113)}, No 2, (2022), 253-262, arxiv/AC:2011.05724v3

 

 \bibitem{P2}  D.\ Popescu, {\em Algebraic valuation ring extensions as limits of complete intersection algebras},  Revista Matematica Complutense,  {\bf 37}, No 2, (2024), 467-472, arxiv/AC:2305.00178.

 
\bibitem{P3}  D.\ Popescu, {\em Pure transcendental, immediate valuation ring extensions as limits of smooth algebras}, to appear in Manuscripta Math., DOI : 10.1007/s00229-025-01635-w,
arxiv/AC:2206.00472v9. 
 

\bibitem{P4}  D.\ Popescu, {\em  Filtered colimits of complete intersection algebras}, to appear in Rev. Roum. Math. Pures Appl., {\bf 70}, No 1-2 (2025),157-165, arxiv/AC:2412.16237v1.

\bibitem{Sch} O.\ F.\ G.\ Schilling, {\em The theory of valuations}, Mathematical Surveys, Number IV, American Math. Soc., (1950).
 \bibitem{S} R.\ Swan, {\em Neron-Popescu desingularization}, in "Algebra and Geometry", Ed. M. Kang, International Press, Cambridge, (1998), 135-192.


\bibitem{Z}  O.\ Zariski, {\em Local uniformization on algebraic varieties}, Ann. of Math. {\bf 41}, (1940), 852-896.  
 




}
\end{thebibliography}

 \end{document}
